\documentclass[12pt,a4paper]{article}
\usepackage[utf8]{inputenc}
\usepackage{ctex}
\usepackage{amsmath}
\usepackage{amssymb}
\usepackage{enumitem}

\title{AI2803-计算机体系结构 作业3}
\author{}
\date{}

\begin{document}

\maketitle

\textbf{一、单选题}

\begin{enumerate}
    \item 以下描述错误的是？（\textbf{B}）
    \begin{itemize}[label={}, leftmargin=*]
        \item A：x86指令长度是不固定的
        \item B：x86指令的操作数必须预存于寄存器中
        \item C：MIPS指令长度是固定的
        \item D：RISC架构的指令种类通常比CISC架构更少
    \end{itemize}

    \item x86体系结构中，寄存器EAX长度为多少位？（\textbf{C}）
    \begin{itemize}[label={}, leftmargin=*]
        \item A：8位 \quad B：16位 \quad C：32位 \quad D：64位
    \end{itemize}

    \item 8086系统中标志位ZF的含义是？（\textbf{B}）
    \begin{itemize}[label={}, leftmargin=*]
        \item A：溢出标志 \quad B：零标志 \quad C：进位标志 \quad 
        \item D：符号标志 \quad E：奇偶标志
    \end{itemize}

    \item 8086系统中段寄存器CS的含义是？（\textbf{A}）
    \begin{itemize}[label={}, leftmargin=*]
        \item A：代码段寄存器 \quad B：数据段寄存器 \quad 
        \item C：附加段寄存器 \quad D：堆栈段寄存器
    \end{itemize}

    \item 8086系统中段寄存器DS的含义是？（\textbf{C}）
    \begin{itemize}[label={}, leftmargin=*]
        \item A：代码段寄存器 \quad B：堆栈段寄存器 \quad 
        \item C：数据段寄存器 \quad D：附加段寄存器
    \end{itemize}

    \item 下列关于MIPS指令的主要特点说法错误的是？（\textbf{C}）
    \begin{itemize}[label={}, leftmargin=*]
        \item A：指令长度固定 \quad B：需要优秀的编译器支持
        \item C：指令数量多，且功能复杂 \quad 
        \item D：只有Load和Store指令可以访问存储器
    \end{itemize}
\end{enumerate}

\textbf{二、多选题}

\begin{enumerate}
    \item IA-32寄存器模型中包括以下哪些寄存器？（\textbf{A, C, D, E}）
    \begin{itemize}[label={}, leftmargin=*]
        \item A：通用寄存器 \quad B：页面寄存器 \quad C：段寄存器 \quad 
        \item D：指令指针寄存器 \quad E：标志寄存器
    \end{itemize}

    \item 下列x86指令中，哪些属于算术运算指令？（\textbf{B, C, G}）
    \begin{itemize}[label={}, leftmargin=*]
        \item A：MOV \quad B：DEC \quad C：MUL \quad D：AND \quad 
        \item E：SHL \quad F：MOVSB \quad G：ADD
    \end{itemize}

    \item MIPS按照指令的基本格式可以分为三种类型，具体为以下哪三种类型？（\textbf{B,C,E}）
    \begin{itemize}[label={}, leftmargin=*]
        \item A：O型指令 \quad B：I型指令 \quad C：J型指令 \quad 
        \item D：M型指令 \quad E：R型指令
    \end{itemize}
\end{enumerate}

\section*{三、问答题}

\begin{enumerate}
    \item 设CS=2500H，DS=2400H，SS=2430H，BX=0200H，SI=0010H，DI=0206H，请分别写出下述x86指令的源操作数的物理地址。
    \begin{enumerate}
        \item MOV AX，[2000H]
        \item MOV AX，[BX+SI+4]
        \item MOV AX，[DI+100H]
    \end{enumerate}

    \textbf{解答：}

    \begin{enumerate}
            \item 26000H
            \item 24214H
            \item 24306H
    \end{enumerate}

    \item 阅读下面的x86汇编程序片段，假设DS=1000H，回答以下问题。
    \begin{verbatim}
    MOV SI, 1250H 
    MOV DI, 1370H 
    MOV CL, 3 
    MOV AX, DS 
    MOV ES, AX 
    MOV BX, 5 
    MOV CH, 1
    STD 
    REP MOVSB
    \end{verbatim}
    (1) 在这次串传送操作中，完成了第一个元素的传送后，SI寄存器的值是什么？ \\
    (2) 在这次串传送操作中，总共传送了多少个字节的数据？

    \textbf{解答：}

    \begin{enumerate}
        \item[(1)] 124FH
        \item[(2)] 259个字节
    \end{enumerate}

    \item 阅读下面的MIPS汇编程序片段，已知第一条指令的地址为0x40000020，第一条指令的16位立即数为0x0008，第三条指令的26位立即数字段为0x0300F0。
    \begin{verbatim}
    beq $t0, $t1, label
    add $t2, $t3, $t4
    j target
    \end{verbatim}
    (1) 对于beq指令，计算分支条件成立时的目标地址（用十六进制表示）；如果分支条件不成立，写出下一条要执行的指令地址（用十六进制表示）。 \\
    (2) 对于j指令，计算跳转的目标地址（用十六进制表示）。 \\
    (3) 如果想跳转至0x30000000，是否可以使用j指令？为什么？
\end{enumerate}

\textbf{解答：}
    \begin{enumerate}
        \item[(1)] 分支条件成立时的目标地址为0x40000044，如果不成立则目标地址为0x40000024。
        \item[(2)] 0x400C03C0
        \item[(3)] 不能。因为j指令的目标地址由当前PC+4的高4位和指令中的26位立即数组成，而当前PC的高四位为0x40，无法直接跳转到0x30开头的地址空间。
    \end{enumerate}
\end{document}